\section{Conclusion}\label{sec:conclusion}

We studied federated LSTM forecasting of household load under a protocol that holds the data, model and budget fixed
and changes one factor at a time, on 100 London homes from the Low Carbon London trial.

Three conclusions follow. First, the federated accuracy gap is a fragmentation cost: the federated protocol on a
single client matches centralised training, and the error grows with the number of clients the same homes are spread
over, to about 10\,\% when every home is its own client. Second, heterogeneity is mainly a communication cost.
Regrouping homes from random to fully profile-clustered barely moved final accuracy but raised the bytes needed to
approach centralised accuracy by about 85\,\%. Third, FL is not the bandwidth-light option here. With a small LSTM and
half-hourly data, federated training moves more bytes than the readings it protects, so its justification is
governance.

For practitioners, this suggests federating a few aggregated clients rather than individual meters, fine-tuning the
global model locally (which removed most of the harm to individual households at community scale), and selecting the
round on validation error. Proximal regularisation brought no benefit. A proximal coefficient scheduled from update
dissimilarity was worse, because that signal is dominated by the stage of convergence. This caution applies to any
regulariser adapted to client-update statistics.

Future work should test these conclusions on other datasets, climates and horizons, and with larger models, where
per-message cost grows. It should also add differential privacy and secure aggregation, whose costs would compound
those measured here, and design heterogeneity signals that are invariant to the stage of training.
